% Options for packages loaded elsewhere
\PassOptionsToPackage{unicode}{hyperref}
\PassOptionsToPackage{hyphens}{url}
%
\documentclass[
]{article}
\usepackage{amsmath,amssymb}
\usepackage{iftex}
\ifPDFTeX
  \usepackage[T1]{fontenc}
  \usepackage[utf8]{inputenc}
  \usepackage{textcomp} % provide euro and other symbols
\else % if luatex or xetex
  \usepackage{unicode-math} % this also loads fontspec
  \defaultfontfeatures{Scale=MatchLowercase}
  \defaultfontfeatures[\rmfamily]{Ligatures=TeX,Scale=1}
\fi
\usepackage{lmodern}
\ifPDFTeX\else
  % xetex/luatex font selection
\fi
% Use upquote if available, for straight quotes in verbatim environments
\IfFileExists{upquote.sty}{\usepackage{upquote}}{}
\IfFileExists{microtype.sty}{% use microtype if available
  \usepackage[]{microtype}
  \UseMicrotypeSet[protrusion]{basicmath} % disable protrusion for tt fonts
}{}
\makeatletter
\@ifundefined{KOMAClassName}{% if non-KOMA class
  \IfFileExists{parskip.sty}{%
    \usepackage{parskip}
  }{% else
    \setlength{\parindent}{0pt}
    \setlength{\parskip}{6pt plus 2pt minus 1pt}}
}{% if KOMA class
  \KOMAoptions{parskip=half}}
\makeatother
\usepackage{xcolor}
\usepackage{color}
\usepackage{fancyvrb}
\newcommand{\VerbBar}{|}
\newcommand{\VERB}{\Verb[commandchars=\\\{\}]}
\DefineVerbatimEnvironment{Highlighting}{Verbatim}{commandchars=\\\{\}}
% Add ',fontsize=\small' for more characters per line
\newenvironment{Shaded}{}{}
\newcommand{\AlertTok}[1]{\textcolor[rgb]{1.00,0.00,0.00}{\textbf{#1}}}
\newcommand{\AnnotationTok}[1]{\textcolor[rgb]{0.38,0.63,0.69}{\textbf{\textit{#1}}}}
\newcommand{\AttributeTok}[1]{\textcolor[rgb]{0.49,0.56,0.16}{#1}}
\newcommand{\BaseNTok}[1]{\textcolor[rgb]{0.25,0.63,0.44}{#1}}
\newcommand{\BuiltInTok}[1]{\textcolor[rgb]{0.00,0.50,0.00}{#1}}
\newcommand{\CharTok}[1]{\textcolor[rgb]{0.25,0.44,0.63}{#1}}
\newcommand{\CommentTok}[1]{\textcolor[rgb]{0.38,0.63,0.69}{\textit{#1}}}
\newcommand{\CommentVarTok}[1]{\textcolor[rgb]{0.38,0.63,0.69}{\textbf{\textit{#1}}}}
\newcommand{\ConstantTok}[1]{\textcolor[rgb]{0.53,0.00,0.00}{#1}}
\newcommand{\ControlFlowTok}[1]{\textcolor[rgb]{0.00,0.44,0.13}{\textbf{#1}}}
\newcommand{\DataTypeTok}[1]{\textcolor[rgb]{0.56,0.13,0.00}{#1}}
\newcommand{\DecValTok}[1]{\textcolor[rgb]{0.25,0.63,0.44}{#1}}
\newcommand{\DocumentationTok}[1]{\textcolor[rgb]{0.73,0.13,0.13}{\textit{#1}}}
\newcommand{\ErrorTok}[1]{\textcolor[rgb]{1.00,0.00,0.00}{\textbf{#1}}}
\newcommand{\ExtensionTok}[1]{#1}
\newcommand{\FloatTok}[1]{\textcolor[rgb]{0.25,0.63,0.44}{#1}}
\newcommand{\FunctionTok}[1]{\textcolor[rgb]{0.02,0.16,0.49}{#1}}
\newcommand{\ImportTok}[1]{\textcolor[rgb]{0.00,0.50,0.00}{\textbf{#1}}}
\newcommand{\InformationTok}[1]{\textcolor[rgb]{0.38,0.63,0.69}{\textbf{\textit{#1}}}}
\newcommand{\KeywordTok}[1]{\textcolor[rgb]{0.00,0.44,0.13}{\textbf{#1}}}
\newcommand{\NormalTok}[1]{#1}
\newcommand{\OperatorTok}[1]{\textcolor[rgb]{0.40,0.40,0.40}{#1}}
\newcommand{\OtherTok}[1]{\textcolor[rgb]{0.00,0.44,0.13}{#1}}
\newcommand{\PreprocessorTok}[1]{\textcolor[rgb]{0.74,0.48,0.00}{#1}}
\newcommand{\RegionMarkerTok}[1]{#1}
\newcommand{\SpecialCharTok}[1]{\textcolor[rgb]{0.25,0.44,0.63}{#1}}
\newcommand{\SpecialStringTok}[1]{\textcolor[rgb]{0.73,0.40,0.53}{#1}}
\newcommand{\StringTok}[1]{\textcolor[rgb]{0.25,0.44,0.63}{#1}}
\newcommand{\VariableTok}[1]{\textcolor[rgb]{0.10,0.09,0.49}{#1}}
\newcommand{\VerbatimStringTok}[1]{\textcolor[rgb]{0.25,0.44,0.63}{#1}}
\newcommand{\WarningTok}[1]{\textcolor[rgb]{0.38,0.63,0.69}{\textbf{\textit{#1}}}}
\usepackage{longtable,booktabs,array}
\usepackage{calc} % for calculating minipage widths
% Correct order of tables after \paragraph or \subparagraph
\usepackage{etoolbox}
\makeatletter
\patchcmd\longtable{\par}{\if@noskipsec\mbox{}\fi\par}{}{}
\makeatother
% Allow footnotes in longtable head/foot
\IfFileExists{footnotehyper.sty}{\usepackage{footnotehyper}}{\usepackage{footnote}}
\makesavenoteenv{longtable}
\setlength{\emergencystretch}{3em} % prevent overfull lines
\providecommand{\tightlist}{%
  \setlength{\itemsep}{0pt}\setlength{\parskip}{0pt}}
\setcounter{secnumdepth}{-\maxdimen} % remove section numbering
\ifLuaTeX
  \usepackage{selnolig}  % disable illegal ligatures
\fi
\IfFileExists{bookmark.sty}{\usepackage{bookmark}}{\usepackage{hyperref}}
\IfFileExists{xurl.sty}{\usepackage{xurl}}{} % add URL line breaks if available
\urlstyle{same}
\hypersetup{
  hidelinks,
  pdfcreator={LaTeX via pandoc}}

\author{}
\date{}

\begin{document}

{
\setcounter{tocdepth}{3}
\tableofcontents
}
\hypertarget{srs-pagetracker}{%
\section{SRS: PageTracker}\label{srs-pagetracker}}

\textbf{Версія:} 1.0 · \textbf{Статус:} узгоджено за SKED-діалогом

\hypertarget{1-ux43fux440ux438ux437ux43dux430ux447ux435ux43dux43dux44f-ux456-ux43cux435ux436ux456}{%
\subsection{1. Призначення і
межі}\label{1-ux43fux440ux438ux437ux43dux430ux447ux435ux43dux43dux44f-ux456-ux43cux435ux436ux456}}

PageTracker --- україномовний вебзастосунок для багатьох незалежних
користувачів. Він дає змогу вести власну бібліотеку книг, фіксувати
послідовно прочитані сторінки й бачити середній темп читання за
календарні періоди. Дані користувача ізольовані від даних інших
користувачів.

До v1 не входять: адміністратори, зміна/відновлення пароля або пошти,
видалення облікового запису, нагадування, інтеграції, імпорт/експорт,
пошук, фільтри, сортування книг, резервне копіювання, юридичні екрани,
спеціальні функції доступності та журналювання технічних помилок.

\hypertarget{2-ux433ux43bux43eux441ux430ux440ux456ux439}{%
\subsection{2.
Глосарій}\label{2-ux433ux43bux43eux441ux430ux440ux456ux439}}

\begin{longtable}[]{@{}ll@{}}
\toprule\noalign{}
Термін & Визначення \\
\midrule\noalign{}
\endhead
\bottomrule\noalign{}
\endlastfoot
Користувач & Особа з окремим обліковим записом і власними даними. \\
Бібліотека & Набір книг одного користувача. \\
Книга & Назва, автор, загальна кількість сторінок, статус і
необов'язковий рейтинг. \\
Запис читання & Діапазон послідовно прочитаних сторінок однієї книги за
поточну дату. \\
\texttt{розпочати} & Статус нової книги без записів. \\
\texttt{читається} & Статус книги з записами, яка ще не завершена. \\
\texttt{прочитано} & Незворотний статус після фіксації останньої
сторінки. \\
Рейтинг & Необов'язкова оцінка завершеної книги від 1 до 5 зірок. \\
\end{longtable}

\hypertarget{3-ux437ux430ux433ux430ux43bux44cux43dux438ux439-ux43eux43fux438ux441}{%
\subsection{3. Загальний
опис}\label{3-ux437ux430ux433ux430ux43bux44cux43dux438ux439-ux43eux43fux438ux441}}

Система має головну сторінку з формою нового запису, бібліотеку,
сторінку деталей книги та сторінку статистики. Після входу відкривається
головна сторінка. Основна навігація: «Головна сторінка», «Бібліотека»,
«Статистика», «Вихід».

\hypertarget{4-ux444ux443ux43dux43aux446ux456ux43eux43dux430ux43bux44cux43dux456-ux432ux438ux43cux43eux433ux438}{%
\subsection{4. Функціональні
вимоги}\label{4-ux444ux443ux43dux43aux446ux456ux43eux43dux430ux43bux44cux43dux456-ux432ux438ux43cux43eux433ux438}}

\hypertarget{41-ux434ux43eux441ux442ux443ux43f}{%
\subsubsection{4.1 Доступ}\label{41-ux434ux43eux441ux442ux443ux43f}}

\begin{itemize}
\tightlist
\item
  \textbf{REQ-F-001 --- Ізоляція даних.} Автентифікований користувач має
  доступ лише до своєї бібліотеки, записів і статистики.
  \textbf{Критерій:} користувач A не може читати або змінювати дані B.
\item
  \textbf{REQ-F-019 --- Реєстрація.} Система створює обліковий запис за
  електронною поштою та непорожнім паролем. Після реєстрації користувач
  виконує окремий вхід.
\item
  \textbf{REQ-F-020 --- Вхід.} Система автентифікує користувача за
  електронною поштою та паролем. \textbf{Критерій:} коректні дані
  відкривають головну сторінку.
\item
  \textbf{REQ-F-021 --- Унікальність пошти.} Пошта унікальна з
  урахуванням регістру: \texttt{User@example.com} і
  \texttt{user@example.com} різні.
\item
  \textbf{REQ-F-022 --- Верифікація пошти.} Підтвердження пошти для
  входу не потрібне.
\item
  \textbf{REQ-F-023 --- Пароль.} Пароль містить щонайменше один символ;
  інших правил складу немає.
\item
  \textbf{REQ-F-029 --- Вихід.} Користувач може завершити сесію; після
  цього дані недоступні без нового входу.
\item
  \textbf{REQ-F-030 --- Помилка доступу.} За зайнятої пошти або
  неправильних даних входу показується «Неправильні дані».
\item
  \textbf{REQ-F-041 --- Після реєстрації.} Система показує форму входу,
  не авторизуючи користувача автоматично.
\item
  \textbf{REQ-F-042 --- Пробіли.} Поля назви, автора, пошти й пароля
  лише з пробілів вважаються порожніми.
\end{itemize}

\hypertarget{42-ux43aux43dux438ux433ux438-ux442ux430-ux440ux435ux439ux442ux438ux43dux433}{%
\subsubsection{4.2 Книги та
рейтинг}\label{42-ux43aux43dux438ux433ux438-ux442ux430-ux440ux435ux439ux442ux438ux43dux433}}

\begin{itemize}
\tightlist
\item
  \textbf{REQ-F-002 --- Створення книги.} Користувач створює власну
  книгу із назвою, автором і загальною кількістю сторінок.
\item
  \textbf{REQ-F-004 --- Статуси.} Книга має статус \texttt{розпочати},
  \texttt{читається} або \texttt{прочитано}; нова книга має
  \texttt{розпочати} і не може одразу мати \texttt{прочитано}.
\item
  \textbf{REQ-F-005 --- Рейтинг.} Нова книга не має рейтингу. Власник
  може встановити або змінити 1--5 зірок лише для \texttt{прочитано}.
\item
  \textbf{REQ-F-013 --- Загальна кількість сторінок.} Її можна змінити
  лише поки книга не має записів.
\item
  \textbf{REQ-F-014 --- Видалення книги.} Перед видаленням показується
  «Ви хочете видалити цю книгу». Після підтвердження видаляються книга й
  усі її записи. \textbf{Критерій:} дані зникають з бібліотеки та
  статистики.
\item
  \textbf{REQ-F-024 --- Метадані.} Назва й автор обов'язкові; дублікати
  назви та автора дозволені.
\item
  \textbf{REQ-F-026 --- Кількість сторінок.} Значення --- ціле число не
  менше 1.
\item
  \textbf{REQ-F-027 --- Некоректна книга.} За порожньої або некоректної
  назви, автора чи кількості сторінок показується «заповніть дані»,
  книга не зберігається.
\item
  \textbf{REQ-F-028 --- Редагування метаданих.} Назву й автора можна
  редагувати за будь-якого статусу.
\end{itemize}

\hypertarget{43-ux437ux430ux43fux438ux441ux438-ux447ux438ux442ux430ux43dux43dux44f}{%
\subsubsection{4.3 Записи
читання}\label{43-ux437ux430ux43fux438ux441ux438-ux447ux438ux442ux430ux43dux43dux44f}}

\begin{itemize}
\tightlist
\item
  \textbf{REQ-F-003 --- Створення запису.} Для однієї книги в одну дату
  дозволено кілька записів. Прочитано сторінок:
  \texttt{кінцева\ −\ початкова\ +\ 1}.
\item
  \textbf{REQ-F-006 --- Завершення.} Запис з останньою сторінкою
  автоматично встановлює \texttt{прочитано}; статус незворотний. Запис
  100--100 для книги на 100 сторінок допустимий.
\item
  \textbf{REQ-F-007 --- Дата.} Система автоматично встановлює поточну
  дату; користувач її не вводить і не редагує.
\item
  \textbf{REQ-F-008 --- Неперетин.} Діапазони однієї книги не
  перетинаються. За перетину зміна не зберігається, показується «Ці
  сторінки вже прочитані».
\item
  \textbf{REQ-F-009 --- Порядок меж.} Якщо кінцева сторінка менша за
  початкову, запис не зберігається, показується «Кінцева сторінка, менша
  за початкову».
\item
  \textbf{REQ-F-010 --- Межі сторінок.} Сторінки --- цілі числа від 1 до
  загальної кількості; для значень поза межами, нецілих або нечислових
  показується «Будь ласка, змініть сторінку».
\item
  \textbf{REQ-F-011 --- Редагування.} Редагувати запис можна лише для
  книги зі статусом \texttt{читається}.
\item
  \textbf{REQ-F-025 --- Початок читання.} Перший незавершальний запис
  змінює статус \texttt{розпочати} на \texttt{читається}.
\item
  \textbf{REQ-F-031 --- Видалення запису.} Видалення доступне лише для
  \texttt{читається}; перед ним показується «Ви хочете видалити цей
  запис».
\item
  \textbf{REQ-F-032 --- Завершення після редагування.} Редагування
  запису до останньої сторінки завершує книгу.
\item
  \textbf{REQ-F-033 --- Повернення статусу.} Видалення останнього запису
  \texttt{читається} повертає статус \texttt{розпочати}.
\item
  \textbf{REQ-F-034 --- Валідація редагування.} Під час редагування
  діють усі правила створення запису.
\item
  \textbf{REQ-F-035 --- Вибір книги.} У формі нового запису доступні
  лише книги \texttt{розпочати} або \texttt{читається}.
\item
  \textbf{REQ-F-036 --- Безперервність.} Перший запис починається зі
  сторінки 1. Кожен наступний починається одразу після попереднього; за
  пропуску показується «Відсутні сторінки» і запис не зберігається.
\item
  \textbf{REQ-F-037 --- Видалення без розриву.} Якщо видалення створить
  розрив у решті записів, система його відхиляє та показує «Видалення
  заборонене».
\item
  \textbf{REQ-F-043 --- Неповна форма.} Якщо не вибрано книгу або не
  заповнено початкову чи кінцеву сторінку, запис не зберігається без
  окремого повідомлення.
\end{itemize}

\hypertarget{44-ux441ux442ux430ux442ux438ux441ux442ux438ux43aux430}{%
\subsubsection{4.4
Статистика}\label{44-ux441ux442ux430ux442ux438ux441ux442ux438ux43aux430}}

\begin{itemize}
\tightlist
\item
  \textbf{REQ-F-015 --- Середнє.} Система обчислює середню кількість
  сторінок за всіма книгами користувача для тижня, місяця або року.
\item
  \textbf{REQ-F-016 --- Періоди.} За замовчуванням показано поточний
  тиждень (понеділок--неділя); доступні минулі тижні, місяці й роки.
\item
  \textbf{REQ-F-017 --- Формула.} За наявності хоча б одного запису
  середнє дорівнює сумі сторінок, поділеній на кількість днів періоду,
  що вже настали. Минулі дні без записів дорівнюють 0, майбутні не
  враховуються.
\item
  \textbf{REQ-F-018 --- Немає даних.} Якщо записів у періоді немає,
  показується «Щоб побачити статистику, додайте новий запис» замість
  числа.
\item
  \textbf{REQ-F-038 --- Формат.} Середнє округлюється до найближчого
  цілого; \texttt{0,5} округлюється вгору. Показується також кількість
  днів у знаменнику.
\item
  \textbf{REQ-F-039 --- Актуальність.} Статистика враховує створення,
  редагування й видалення записів.
\item
  \textbf{REQ-F-040 --- Технічна помилка.} За технічної помилки
  показується «Помилка».
\end{itemize}

\hypertarget{5-ux432ux438ux43cux43eux433ux438-ux434ux43e-ux456ux43dux442ux435ux440ux444ux435ux439ux441ux443}{%
\subsection{5. Вимоги до
інтерфейсу}\label{5-ux432ux438ux43cux43eux433ux438-ux434ux43e-ux456ux43dux442ux435ux440ux444ux435ux439ux441ux443}}

\begin{itemize}
\tightlist
\item
  \textbf{REQ-I-001.} Система є вебзастосунком у браузері.
\item
  \textbf{REQ-I-002.} Усі елементи та повідомлення інтерфейсу
  українською.
\item
  \textbf{REQ-I-003.} Сторінка деталей містить метадані, статус,
  рейтинг, записи й доступні дії.
\item
  \textbf{REQ-I-004.} Бібліотека показує назву, автора, статус,
  кількість сторінок і рейтинг за наявності.
\item
  \textbf{REQ-I-005.} Окрема сторінка статистики показує період, середнє
  і кількість днів або повідомлення про відсутність даних.
\item
  \textbf{REQ-I-006.} У списку записів показані лише дата і кількість
  сторінок; порядок --- від найстарішої дати до найновішої, у межах дати
  --- за зростанням початкової сторінки.
\item
  \textbf{REQ-I-007.} Запис створюється на головній сторінці.
\item
  \textbf{REQ-I-008.} Форма запису містить вибір книги, початкову та
  кінцеву сторінки; дата й кількість не вводяться.
\item
  \textbf{REQ-I-009.} Після входу відкривається головна сторінка.
\item
  \textbf{REQ-I-010.} Книга створюється формою у бібліотеці.
\item
  \textbf{REQ-I-011.} Після створення відкривається бібліотека з новою
  книгою.
\item
  \textbf{REQ-I-012.} Доступні розділи «Головна сторінка», «Бібліотека»,
  «Статистика» та «Вихід».
\item
  \textbf{REQ-I-013.} Запис редагується на сторінці деталей відповідної
  книги.
\item
  \textbf{REQ-I-014.} На сторінках входу й реєстрації є переходи між
  ними.
\item
  \textbf{REQ-I-015.} Бібліотека поділена на сторінки по 16 книг.
\item
  \textbf{REQ-I-016.} Перехід між сторінками здійснюється номерами
  сторінок.
\item
  \textbf{REQ-I-017.} Після видалення книги відкривається бібліотека.
\item
  \textbf{REQ-I-018.} Метадані книги редагуються на сторінці її деталей.
\item
  \textbf{REQ-I-019.} Після створення запису відкривається головна
  сторінка.
\end{itemize}

\hypertarget{6-ux43dux435ux444ux443ux43dux43aux446ux456ux43eux43dux430ux43bux44cux43dux456-ux432ux438ux43cux43eux433ux438}{%
\subsection{6. Нефункціональні
вимоги}\label{6-ux43dux435ux444ux443ux43dux43aux446ux456ux43eux43dux430ux43bux44cux43dux456-ux432ux438ux43cux43eux433ux438}}

\begin{itemize}
\tightlist
\item
  \textbf{REQ-NF-001 --- Сумісність.} Система коректно працює в сучасних
  настільних браузерах.
\item
  \textbf{REQ-NF-002 --- Продуктивність.} За активності до 10
  користувачів вхід, відкриття бібліотеки, створення запису й показ
  статистики виконуються за ≤ 2 секунди.
\item
  \textbf{REQ-NF-003 --- Безпека.} Паролі не зберігаються відкрито,
  передавання даних відбувається через HTTPS.
\item
  \textbf{REQ-NF-004 --- Збереженість.} Збережені дані доступні після
  оновлення сторінки, виходу й повторного входу.
\item
  \textbf{REQ-NF-005 --- Сесія.} Сесія не завершується автоматично, доки
  користувач не натисне «Вихід».
\item
  \textbf{REQ-NF-006 --- Атомарність.} За технічної помилки створення
  або редагування книги чи запису дані не зберігаються частково.
\item
  \textbf{REQ-NF-007 --- Місткість.} Підтримуються до 100 книг і 5 000
  записів на користувача.
\end{itemize}

\hypertarget{7-ux43fux440ux43eux454ux43aux442ux43dux456-ux43eux431ux43cux435ux436ux435ux43dux43dux44f}{%
\subsection{7. Проєктні
обмеження}\label{7-ux43fux440ux43eux454ux43aux442ux43dux456-ux43eux431ux43cux435ux436ux435ux43dux43dux44f}}

\begin{itemize}
\tightlist
\item
  Серверна частина: \textbf{Node.js}; клієнтська: \textbf{React}; база
  даних: \textbf{PostgreSQL}.
\item
  Часовий пояс: \texttt{Europe/Kyiv}.
\item
  У v1 немає пошуку, фільтрів і сортування книг; записи відображаються
  одним списком без сторінок.
\item
  Формат електронної пошти не перевіряється; пошта не верифікується;
  кількість невдалих входів не обмежується.
\item
  Не встановлюються SLA, резервне копіювання, спеціальна доступність,
  юридичні екрани та журналювання технічних помилок.
\end{itemize}

\hypertarget{8-ux43aux43bux44eux447ux43eux432ux456-ux441ux446ux435ux43dux430ux440ux456ux457-ux43fux440ux438ux439ux43dux44fux442ux442ux44f}{%
\subsection{8. Ключові сценарії
прийняття}\label{8-ux43aux43bux44eux447ux43eux432ux456-ux441ux446ux435ux43dux430ux440ux456ux457-ux43fux440ux438ux439ux43dux44fux442ux442ux44f}}

\begin{Shaded}
\begin{Highlighting}[]
\NormalTok{Scenario: Завершення книги}
\NormalTok{  Given книга "читається" і має 100 сторінок}
\NormalTok{  When користувач зберігає допустимий запис 91–100}
\NormalTok{  Then статус стає "прочитано"}
\NormalTok{  And її записи не можна редагувати або видаляти}

\NormalTok{Scenario: Пропуск сторінок}
\NormalTok{  Given існує запис 1–10}
\NormalTok{  When користувач створює запис 12–20}
\NormalTok{  Then система не зберігає запис}
\NormalTok{  And показує "Відсутні сторінки"}

\NormalTok{Scenario: Видалення проміжного запису}
\NormalTok{  Given книга має записи 1–10, 11–20, 21–30}
\NormalTok{  When користувач підтверджує видалення 11–20}
\NormalTok{  Then система не видаляє запис}
\NormalTok{  And показує "Видалення заборонене"}

\NormalTok{Scenario: Немає статистики}
\NormalTok{  Given у вибраному періоді немає записів}
\NormalTok{  When користувач відкриває статистику}
\NormalTok{  Then система показує "Щоб побачити статистику, додайте новий запис"}
\end{Highlighting}
\end{Shaded}

\hypertarget{9-ux43aux43eux43dux432ux435ux43dux446ux456ux44f-ux43aux43eux43cux456ux442ux456ux432}{%
\subsection{9. Конвенція
комітів}\label{9-ux43aux43eux43dux432ux435ux43dux446ux456ux44f-ux43aux43eux43cux456ux442ux456ux432}}

Формат: \texttt{префікс:\ стислий\ опис}.

\begin{itemize}
\tightlist
\item
  \texttt{spec:} --- специфікації в \texttt{/spec}; \texttt{test:} ---
  тести; \texttt{feat:} --- функціональність; \texttt{docs:} ---
  документація; \texttt{gate:} --- релізи та CI/CD.
\end{itemize}

\end{document}
